\documentclass[../../../include/open-logic-section]{subfiles}

\begin{document}

\olfileid{sth}{cardinals}{zfc}
\olsection{$\ZFC$：一个里程碑}

随着良序公理的加入，我们到达了理论上的最后一个里程碑。我们现在已具备构成~$\ZFC$ 所需的全部公理。具体来说：

\begin{defn}
理论 $\ZFC$ 包含以下公理：外延公理、并集公理、配对公理、幂集公理、无穷公理、基础公理、良序公理，以及分离公理模式和替换公理模式的所有实例。换言之，$\ZFC$ 就是在~$\ZF$ 的基础上加入良序公理。
\end{defn}

$\ZFC$ 代表带\emph{选择公理}的 \emph{Zermelo--Fraenkel}集合论。这可能看似有些奇怪，因为我们添加的公理称为“良序公理”，而不是“选择公理”。但是，当我们后续阐述选择公理时，将会发现良序公理与选择公理（在~$\ZF$ 意义下）是等价的（参见 \olref[choice][woproblem]{thmwochoice}）。因此，将哪一个作为我们的“基本”公理都无区别。而且，“$\ZFC$”是文献中的标准名称。

\end{document}